\section{Discussion and Limitations}
The current evidence supports distinct claims. Repeated small benchmarks evaluate formulation checking;
the synthetic suite tests workflow connections; frozen models test solver behaviour; specified-source pilots
test particular data and modelling paths. None substitutes for repeated autonomous real-data acceptance.
The offline suite passes 146 tests, including schema, checker, recovery and checkpoint fixtures, but test counts
do not measure semantic correctness or live acquisition reliability.

Runtime checkers can share LLM and schema errors. Matching counts and totals can hide per-entity mistakes,
and sampled requirement checks cannot verify every state. Exhaustive ground-state checking is limited to 26
variables. A generated objective can still replace the requested objective with a weaker surrogate.
Model repair and data substitution may require user judgment, even when structural checks pass.

Hardware comparisons contain three repetitions per condition and allocate different service budgets.
They establish working submissions and instance-specific outcomes, not quantum advantage, lower total cost or
a generally preferable backend. Conventional baselines can be the delivered result, as in the network pilot.

\section{Conclusion}
We presented a local-LLM workflow connecting requirements, typed data preparation, modelling, annealing and
checked delivery. Separate validation of data, model meaning and solver feasibility supports bounded repair.
Ground-state verification exposes reference-bitstring failures. The next step is a frozen autonomous real-data
comparison measuring the incremental effects of checking, repair and retrieval.
